\begin{textsample}{POS Dim 7 – global\_south\_2025\_03 – Score 4.00 – global\_south\_2025\_03/122542802.txt}  \label{ex:f7_pos_427}
The United Nations : A Toothless Tiger?
The United Nations Security Council, crippled by power politics and veto-wielding giants, has failed to protect civilians in Gaza, Ukraine, and beyond -- raising urgent questions about its legitimacy and the need for reform.
Israel always intended to continue its relentless pounding of the people in the enclave in its attempt to stamp out the Hamas.
Will the United Nations, tasked with maintaining peace and security in the world, restrain Israel?
It will not, as Israel is backed by the US.
The United Nations Security Council ( UNSC ), the highest decision making body in the UN, has failed to protect civilians in Gaza because of America ' s iron-clad support for the Jewish state.
But it is not just in Gaza that the UNSC has failed.
It has done so in Ukraine, and Syria, as well as in the ongoing civil wars across Africa, including in Sudan and the \textbf{Democratic} Republic of Congo.
The UNSC has @ @ @ @ @ @ @ @ @ @ China, Russia, France, and UK--to use their veto powers to protect friends and punish enemies.
The UNSC no longer reflects the world as it is today.
Neither Britain nor France are as powerful as they were in the past.
It was formed in 1945, when the victors of World War II decided to take charge of the world.
Today, the UNSC has lost all relevance and is now a toothless tiger.
The last meaningful reform was in 1965 when it expanded from 11 members to 15, adding four \textbf{seats} to the two-year elected membership.
Non-permanent members do not have veto rights."The international organisations that were created have become almost irrelevant, there is no reform in them.
Institutions like the UN can not play their role.
People in the world who do not care about laws and rules are doing everything, no one is able to stop them,"Prime Minister Narendra Modi said on Lex Fridman ' s podcast on Sunday.
Modi is not the @ @ @ @ @ @ @ @ @ @ Council has ceased to be the guarantor of world security and has become a battleground for the political strategies of only five countries,"President Recep Tayyib Erdogan of Turkiye said recently.
UN Secretary General, Antonio Guterres had warned that the stark choice before the UN is"to reform or bust ' '."The world has changed.
Our institutions have not,"Guterres said in an address to the General Assembly earlier."We can not effectively address problems as they are if institutions do not reflect the world as it is.
Instead of solving problems, they risk becoming part of the problem,"he said.
Guterres and the UN staff have been crying themselves hoarse over the treatment of civilians in Gaza and affected civilians in strife torn regions of the world.
His is a strong voice, but the UN has become as helpless as a non-profit organisation that can speak out but relies on the governments to take action.
Soon after the UNSC was formed, @ @ @ @ @ @ @ @ @ @, the US and the Soviet Union, against each other.
Both countries used their veto powers at the UNSC to ensure that any resolution that went against their self-interest was defeated.
This has continued to this day.
For a short while, after the break up of the Soviet Union, the UNSC tried to play a positive role ; but it soon got bogged down by rivalries of the P5 members.
Israel ' s disregard for international covenants of war has much to do with the fact that Prime Minister Benjamin Netanyahu has strong American backing.
Every resolution at the UNSC against Israel is vetoed by the US.
Similarly, China would constantly block resolutions on North Korea, Taiwan, Myanmar or Pakistan.
Russia ensured that anti-Iran resolutions were vetoed. \textbf{Republicans} in the US have been critical of the UN, regarding the organisation as a white elephant.
George Bush bypassed the UNSC when he ordered the invasion of Iraq in 2003.
Getting the approval of the UNSC gives military action a measure @ @ @ @ @ @ @ @ @ @ Trump and team have always shown their dislike for the UN and its bodies.
Trump has pulled out of the Paris climate agreement as well as from the World Health Organisation.
UNSC Reforms During his tenure as UN Secretary General, Kofi Annan pushed for expansion of the UNSC, saying the P5 had to reflect the changed world.
It is nearly two decades since a serious push for expansion began, but has so far yielded little result.
India believes it has the strongest claim as one of the most populated countries in the world.
Nearly all countries, from the US to Russia and several EU countries, have promised support for India ' s membership.
Japan is another aspirant from Asia.
Germany is a strong \textbf{candidate} from Europe.
Both Japan and Germany contribute liberally to the UN.
Italy believes it too should sit at the high table.
Brazil too is a contender as was Argentina once.
South Africa and Egypt both think they should represent the African continent, as @ @ @ @ @ @ @ @ @ @, Brazil, Germany and Japan are working together as the Group of 4, to reform the UN.
In September 2024, US announced its support for two permanent \textbf{seats} in an expanded UNSC for Africa.
The Biden administration, keen to repair ties with Africa to counter China ' s spread across the continent, announced this sop to woo the continent.
However, so far there is no clarity or consensus on how many more members will be part of the reformed UNSC.
None of the P5 countries are so far willing to grant veto rights to the new entrants, whenever that happens.
Any reform of the Security Council would require the agreement of all five, and at least two-thirds of UN member states.
That is a tough ask.
But the larger question remains if a reformed UNSC will also be in a position to work cooperatively on maintaining peace and security in the world.

% matched lemmas: candidate, democratic, republican, seat
\end{textsample}
